%% kapa: exam=1 round=24 year=2013 period=1 form=A subject=civil_law subject_ko=민법 seq=35 number=35
%% source: Q-Net 감정평가사 기출문제 (cst00309) artlSeq=5073802 file=extracted/1008/5073802/2013년 제24회 감정평가사 1차시험 기출문제/1교시(A형).hwp
%% conversion: kordoc 4.24.51 -> tools/gen1.py (auto); review=proofread-sonnet; flags=none
\begin{question}{35}
물상보증에 관한 설명으로 옳지 않은 것은? (다툼이 있으면 판례에 의함)
\begin{choices}
  \item 채무자의 채무를 변제한 물상보증인은 변제자대위에 의하여 채권자의 권리를 행사할 수 있다.
  \item 물상보증인이 제공한 담보부동산의 제3취득자가 피담보채무의 이행을 인수하더라도 저당권이 실행되면 물상보증인만이 채무자에게 구상권을 행사할 수 있다.
  \item 물상보증인의 채무자에 대한 구상권에는 민법상 일반채권의 소멸시효가 적용된다.
  \item 근저당권을 설정한 물상보증인은 채권최고액을 변제하여 근저당권설정등기의 말소를 청구할 수 있다.
  \item 채권자가 물상보증인이 제공한 담보목적물에 대하여 근저당권 실행 경매를 신청한 경우 채무자에게 경매개시결정이 송달되지 않더라도 피담보채권의 소멸시효가 중단된다.
\end{choices}
\end{question}
\answer{5}
